\section{\systemname Runtime}
\label{sec:method}

\subsection{System Overview}

\systemname is a verification-centered runtime built around a user, four
model-driven roles, and three system planes. The user supplies standing intent,
domain judgment, and authority over material tradeoffs. The \emph{control plane}
turns that intent into a durable task contract and schedules work; the
\emph{execution plane} performs one bounded mission against real tools and
artifacts; and the \emph{record plane} stores evidence, decisions, and contract
history without deciding whether the work is scientifically complete. The
distinction is operational: control owns scheduling and contract transitions,
execution owns work and mission review, and the record plane owns the immutable
record but no completion decision.

The four roles divide research authority as summarized in
Table~\ref{tab:roles}. The Manager spans the campaign rather than one execution
round. The Manager is also the user-facing boundary: it preserves standing intent,
surfaces consequential ambiguity, and commits authorized contract refinements. The
Planner authors future work and the Engineer performs it. Every normal mission
round then passes to a fresh independent Reviewer; historical
Engineer-self-review fields remain only for compatibility with older traces. The
interfaces are structured objects rather than untyped prose, which makes ownership
and recovery behavior explicit.

The notation $K_t$ is a report-level abstraction over persisted campaign identity,
objectives, Stage state, constraints, verifier obligations, and decision records.
The current implementation distributes these fields across runtime surfaces rather
than storing one monolithic object named $K_t$.

\begin{table}[H]
\centering
\small
\begin{tabularx}{\linewidth}{@{}p{1.6cm} p{4.0cm} X@{}}
\toprule
\rowcolor{papersand}
\textbf{Role} & \textbf{Primary responsibility} & \textbf{Authoritative output} \\
\midrule
Manager & Preserve standing intent; commit task contract, lifetime, route, and campaign stage & Contract refinement, user escalation, campaign division, and stage transition \\
Planner & Decompose the current verified frontier into bounded tasks and dependencies & Task specifications, plan status, verifier obligations, and stage checklist updates \\
Engineer & Modify artifacts, invoke tools, and run experiments for one round & Execution result, artifacts, measurements, and checkpoint handoff \\
Reviewer & Independently inspect every normal round against the current contract & \code{done}, \code{continue}, \code{blocked}, or \code{replan\_requested}; verified frontier update; contradiction or expert-question report \\
\bottomrule
\end{tabularx}
\caption{Authority separation in \systemname. The Engineer executes; a fresh
independent Reviewer controls admission after every normal mission round; Manager
owns contract and Stage transitions at the user-facing boundary.}
\label{tab:roles}
\end{table}

\begin{paperbox}{Algorithm 1: Reviewed mission loop}
\small
\begin{tabbing}
\quad\=\quad\=\quad\=\kill
\textbf{Input:} standing intent $\iota$, contract $K_t$, user state $X_t$, runtime state $H_t$, backlog $B$, budget $b$ \\
$c \leftarrow \operatorname{ManagerCommit}(\iota,K_t)$; persist campaign identity and contract \\
\textbf{while} $b$ remains and $c$ is not complete \textbf{do} \\
\> $q \leftarrow \operatorname{Planner}(B,H_t,c)$; claim one bounded task \\
\> $\tau_t \leftarrow [\,]$ \\
\> \textbf{repeat} \\
\>\> $(y,e) \leftarrow \operatorname{Engineer}(q,H_t)$ \\
\>\> $r \leftarrow \operatorname{Reviewer}(K_t,q,y,e)$ \\
\>\> append $(q,y,e,r)$ to $\tau_t$ \\
\>\> \textbf{if} $r=\code{continue}$ \textbf{then} $q \leftarrow r.\code{next\_action}$ \\
\> \textbf{until} $r\in\{\code{done},\code{blocked},\code{paused},\code{replan\_requested}\}$ \\
\> $E_t \leftarrow \operatorname{AdmitResult}(e,r)$ \\
\> $X_{t+1} \leftarrow \operatorname{SurfaceEvidence}(X_t,E_t,r)$ \\
\> $K_{t+1} \leftarrow K_t$ \\
\> \textbf{if} $r=\code{replan\_requested}$ or contract refinement is proposed \textbf{then} \\
\>\> $K_{t+1} \leftarrow \operatorname{ManagerAdmit}(K_t,E_t,r,u_t)$ \\
\> $H_{t+1} \leftarrow U(H_t,\tau_t,E_t,K_{t+1})$; persist outcome, trace, contract, and state \\
\> $t \leftarrow t+1$; refill $B$ when required \\
\textbf{end while} \\
\textbf{Output:} accepted frontier, inspectable trajectory, refined task contract, and updated expert/runtime state
\end{tabbing}
\end{paperbox}

\subsection{Bounded Missions over Durable Contracts}

The long-lived object in \systemname is the campaign, not a provider transcript. A
campaign has a persisted identity, standing intent, current task contract, accepted
frontier, and explicit unresolved questions. Its work is divided into missions with
explicit outcomes. The scheduler assigns one mission at a time, records its result,
and advances only at a clean mission boundary. Mission assignment is transactional,
preventing duplicate work under concurrency, and resumed execution is tied to the
persistent campaign identity.

Engineer and Reviewer calls use fresh provider sessions for each round. Cross-session
continuity comes from an ordinary shared \code{CHECKPOINT.md} containing durable
state, evidence references, open questions, and the next step. This follows the
broader move from monolithic context toward managed memory tiers and reusable
workflows~\citep{packer2023memgpt,wang2024workflowmemory,xu2025amem}. The Engineer
updates the checkpoint after execution. The Reviewer then reads and corrects the
same file and is its final editor for that round. The full history remains in
artifacts and the event record.

This design makes restart and upgrade behavior part of the method. A running process
hands off only at a mission boundary, and the replacement resumes from the same
persistent campaign identity. After an interruption, replay or reassignment begins
from committed campaign state rather than reconstructing the task from a model
transcript.

The same persistence supports user interaction at a larger timescale. The user
does not need to repeat the original context or monitor every tool call. When a
Reviewer finds that the current verifier cannot distinguish alternatives, that a
constraint conflicts with observed reality, or that multiple valid tradeoffs remain,
the issue is surfaced with evidence and options. The resulting clarification is
committed through Manager into the durable contract. Later missions therefore begin
from the refined understanding rather than from either the original ambiguous prompt
or an agent-invented replacement goal.

\subsection{Autonomy between User Decision Boundaries}

The runtime does not require a human to supervise each round. The harness assigns
local authority to the four roles: Manager maintains the campaign and contract,
Planner creates bounded executable work, Engineer acts on real artifacts, and an
independent Reviewer decides whether evidence supports completion, continuation,
blocking, or replanning. Checkpoints, event history, task claims, external-job
state, and mission outcomes survive session replacement. Consequently, the runtime
can execute for long intervals without synchronous human input and surface a user
decision only when evidence exposes a material value tradeoff, missing information,
or operator-owned constraint.

This is stronger than adding a human-approval button. The autonomous path must
remain capable of selecting routes, running experiments, rejecting invalid work,
consolidating intermediate progress, and ending an unproductive mission. Human
input refines the campaign at consequential boundaries; it is not the mechanism
that keeps the ordinary loop alive.

\subsection{Verification as the Pivot Gate}

All role invocations pass through a common instrumentation layer that records usage
and associates each call with the mission trace. The Engineer and Reviewer share
the same artifact state, allowing review of the actual outputs and execution record
rather than only the Engineer's summary. A typed, append-only trace is the canonical
timeline; user interfaces are projections over that record.

The source of a mission completion verdict is explicit: a fresh Reviewer returns a
structured verdict after every normal Engineer round. Review evaluates more than
terminal completeness. It can admit a bounded frontier increment, reject an
unsupported claim, identify a faulty verifier or inconsistent specification,
request another experiment, or ask Manager to replan and surface a user decision.
Historical
\code{review\_source=engineer\_self\_review} values describe earlier runtime
versions and are retained only so longitudinal traces remain interpretable.
Completion never comes from filenames, Token volume, content hashes, or prose
keywords. The corresponding artifacts and execution records remain inspectable,
while credential redaction preserves the evidence boundary.

Verification is therefore the gate between persistence and pivoting. Agent
execution externalizes hypotheses into artifacts and measurements; Reviewer turns
those outputs into calibrated claims and unresolved alternatives; Manager presents
material choices in terms the user can inspect; and user judgment updates the
contract governing subsequent execution. Without verification, persistence can
compound error. Without persistent self-evolution, verification repeatedly finds
the same errors. Their combination converts a long trajectory into cumulative
problem understanding and reusable capability.

\subsection{Staged Approach Self-Evolution and Contract Pivoting}

Equation~\ref{eq:runtime-update} separates the model from the state around it.
Table~\ref{tab:runtime-ownership} identifies how each component can change and who
owns the committed update. The ownership model is intentionally not uniform. Memory
and skills use a work-versus-certification split: the Engineer or Scientist produces
a candidate, and the Reviewer commits the retained form. Tools and procedures are
system-configured. Stage checklists are Planner-owned, with Reviewer feedback rather
than a second commit gate. Routing policy is Manager-committed, while task
definitions are Planner-authored and scheduler-committed. The task contract is
different from all of these: evidence and Reviewer findings may propose a change,
Manager commits operational refinements, and the user retains authority
over standing intent and value-laden tradeoffs that evidence alone cannot settle.

We use verifier-gated runtime self-evolution to mean the staged evolution of the approach
used to move from the current objective toward satisfaction of the current
constraints. The underlying model parameters remain fixed, while reviewed missions
change the persistent search state available to later work. Each Stage maintains an
approach record $\Pi_t$ containing route hypotheses, actions, task-native evidence,
Reviewer verdicts, retained failures, and the next admissible frontier. A complete
update cycle has four parts:
(1) an execution trajectory produces a candidate memory, skill, procedure,
verification rule, routing decision, or task definition; (2) the responsible role
checks the candidate against artifacts and task-native evidence; (3) the authorized
owner commits, revises, or rejects the update; and (4) a later mission retrieves the
retained state as part of its starting context or execution policy. Activity that
does not survive this commit-and-reuse path is not counted as self-evolution.
Contract pivoting is distinct: it is the evidence-backed refinement of $K_t$ and
$X_t$, with user clarification or authorization where required. Neither mechanism implies that every
mission improves the system or changes the contract.

This preserves a strong self-evolution claim without conflating it with goal
mutation. $K_t$ defines the current objective, constraints, and verifier;
$\Pi_t$ records how the runtime has learned to navigate inside that envelope. If a
verified contract refinement changes the envelope, the prior approach record
remains staged evidence of which routes failed and which procedures or verifiers
remain reusable.

\begin{table}[H]
\centering
\small
\begin{tabularx}{\linewidth}{@{}p{2.25cm} p{2.75cm} p{2.55cm} X@{}}
\toprule
\rowcolor{papersand}
\textbf{State} & \textbf{Change source} & \textbf{Commit owner} & \textbf{Persistent form} \\
\midrule
$M$: memory & Engineer trajectory & Reviewer & Curated checkpoint and event-derived journal \\
$S$: skills & Engineer / Scientist & Reviewer & Versioned skill library \\
$A$: tools and procedures & System configuration & System configuration & Versioned tool and procedure registry \\
$V$: verification & Planner & Planner & Stage checklist; Reviewer supplies feedback \\
$R$: routing and roles & Runtime policy & Manager & Campaign routing policy \\
$Q$: tasks and evaluations & Planner & Scheduler & Backlog task specifications and scopes \\
$\Pi$: staged approach & Mission trajectories and Reviewer verdicts & Manager / Planner & Accepted frontier, rejected routes, route rationale, and next approach \\
$K$: task contract & User clarification, evidence, and Reviewer findings & Manager / user & Versioned intent, objective, constraints, verifiers, and provenance \\
\bottomrule
\end{tabularx}
\caption{Attributable ownership of persistent and revisable state. A mission usually updates only
a subset of these components; material changes to standing user intent are not
delegated to the runtime.}
\label{tab:runtime-ownership}
\end{table}

Two knowledge surfaces carry reusable experience. Versioned skills store procedures
that can be matched to later tasks. A project knowledge base stores source-linked
records and synthesized pages derived from reviewed outcomes. Neither surface is
treated as automatically correct: entries can be revised, archived, or retired when
later results contradict them.

This mechanism can improve a later mission without changing the model itself. A
retained failure can prevent a repeated dead end; a reviewed procedure can shorten
environment inspection; a verifier can reject a previously accepted shortcut; and
a routing update can assign independent review to a higher-risk task. The mechanism
does not imply monotonic improvement. Some missions commit no reusable state,
retained state can become stale, and a harder task distribution can increase cost
even after useful state has accumulated.

The user-facing benefit is complementary. A retained failure can show that a
requested optimization is physically or economically irrelevant; a stronger
verifier can expose that the original success metric was gameable; and a proof
counterexample can replace a vague conjecture with a sharper reachable frontier.
The runtime is not learning the user's preferences covertly. It is making the
evidence and tradeoffs explicit enough for those preferences to be stated and
revised deliberately.

\subsection{Reliability and Resource Governance}

Long-running systems fail through operational drift even when individual model
calls are strong. \systemname therefore combines round-level progress
classification with mission-boundary replanning. Work that repeatedly produces no
decision or artifact change is stopped or reformulated, while long-running external
jobs are tracked separately from model reasoning. These mechanisms bound execution;
scientific correctness remains with the task evaluator and independent Reviewer.

Short Planner DAG nodes expose a subtle control problem: their default three-round
budget is much smaller than the default absolute thresholds for semantic stall,
soft escalation, and hard escalation. The current runtime fits each enabled guard
to the reachable range when \code{max\_rounds} shrinks. For budget $m$, semantic
stall and soft escalation must satisfy $g\leq m-1$, while hard escalation may fire
on the final round and therefore satisfies $g\leq m$. Explicit zero still disables
a guard, and the default 500-round configuration is unchanged. A two-round budget
does not silently turn a configured multi-round stall detector into a one-strike
policy: the impossible detector is disabled unless the caller explicitly requests
threshold one. This coordination preserves meaningful stop and replan signals
instead of falling through silently to the generic maximum-round outcome.

These liveness guards do not decide scientific truth and do not freeze a mistaken
contract. They create bounded points at which Reviewer and Planner can distinguish
three cases: the current route is unproductive, the evidence frontier has advanced
but terminal work remains, or the task contract itself needs user-visible
refinement. This distinction is important because ``not done'' is not synonymous
with ``no progress,'' and ``cannot verify'' is not synonymous with ``incorrect.''

We call this property \emph{honest stopping}: the harness can end the current unit
of work without converting uncertainty, exhaustion, or an external dependency into
a success claim. The implementation keeps completion, incomplete research,
stalling, blocking, infrastructure blocking, failure, and operator interruption as
separate outcome dimensions. Reviewer may return \code{blocked} or
\code{replan\_requested}; a hard escalation ends the mission with a
planner-readable reason. Semantic stall accumulates only from the Reviewer's
explicit structured \code{forward\_progress=false}; a missing or malformed signal
is treated as unknown rather than inferred from prose or filesystem activity.
Stopping therefore preserves both epistemic honesty and the information needed to
resume, decompose, or change approach.

Resource accounting is centralized at model-call and external-job boundaries.
Budgets are checked before work begins and reconciled after completion, preventing
individual roles from expanding their own allocation. Concrete scheduling,
sandboxing, and deployment mechanisms are implementation choices rather than part
of the scientific claim.
